\section{Multi Target Tracking}

% \begin{figure}
%     \incfig[]{tracker}
% \end{figure}

% \begin{itemize}
%     \item Multi-target tracking converts a sequence of noisy radar detections into stable target tracks over time.
%     \item Its purpose is to estimate the current state of each target, including position and velocity, while maintaining continuity between successive radar updates.
%     \item In radar systems, tracking is typically performed after detection and angle estimation, using range, Doppler, and angle measurements as inputs.

%     \item A tracking filter is used to smooth noisy measurements and predict target motion between radar updates.
%     \item Kalman filtering provides a common framework for this process by combining a motion model with uncertain measurements to produce an improved estimate of target state.
%     \item The filter alternates between prediction and update stages, using the predicted target state to guide association and the new measurements to refine the estimate.
%     \item Variants such as the Extended Kalman Filter may be required when the measurement model is nonlinear.

%     \item Radar measurements are naturally produced in polar coordinates, typically as range, bearing, elevation, and sometimes Doppler.
%     \item For tracking, these measurements may be processed either directly in polar form or transformed into Cartesian coordinates depending on the filter design and motion model.
%     \item The choice of coordinate system affects the linearity of the model, the form of the measurement uncertainty, and the complexity of the update equations.

%     \item Gating is used to restrict which detections are considered plausible matches for an existing track.
%     \item This reduces computational cost and helps reject unlikely associations caused by clutter or nearby targets.
%     \item Data association then determines which measurement should be assigned to each predicted track when multiple candidates are present.

%     \item Track life management controls how tracks are initiated, confirmed, maintained, and deleted over time.
%     \item New detections may first form tentative tracks, which are only promoted to confirmed tracks after consistent observations across multiple updates.
%     \item Similarly, tracks may be deleted after repeated missed detections to avoid maintaining false or stale targets.
%     \item These stages allow the tracker to balance responsiveness to new targets with robustness against clutter and intermittent measurement loss.
% \end{itemize}

Multi-target tracking (MTT) converts a sequence of noisy radar detections into stable 
estimates of target state over time. Each radar update produces a set of detections 
characterised by range, bearing, and elevation, as described in the preceding sections. 
The role of the tracker is to associate these detections with targets across 
successive updates, filter out noise, and maintain a continuous estimate of each 
target's state to support detect-and-avoid (DAA) decision making.

A tracking filter combines a model of target motion called the transition model with 
incoming measurements to produce an evolving state estimate over time. The Kalman 
filter \cite{jazwinski1970stochastic} is 
a standard framework, alternating between a prediction stage, in which the current 
state is propagated forward using a motion model, and an update stage (see Figure 
\ref{fig:kalman_filter}), in which a new 
measurement corrects the prediction. A constant velocity (CV) with noise is a 
common motion model, and is represented in three dimensions with a matrix for 
the state transition and a covariance matrix for the process noise. For manoeuvring 
targets the interacting multiple model (IMM) \cite{blom1988imm} 
maintains a bank of filters under different motion assumptions and blends their outputs 
probabilistically.

\begin{figure}[htb]
    \incfig[0.6]{kalman_filter}
    \label{fig:kalman_filter}
    \caption{Each time step, the makes a prediction of the target's current state based
    on the previous state and a motion model (red), and then updates the prediction based 
    on the new measurement (yellow) to produce an improved estimate (blue). The detection, 
    prediction and updated estimation are each represented as a Gaussian distribution, 
    hence the ellipses, with the mean representing the most likely state and the 
    covariance representing the uncertainty.}
\end{figure}


Since radar measurements are expressed in spherical coordinates while motion models are 
most naturally formulated in Cartesian space, the measurement equation is nonlinear. 
The Extended Kalman Filter (EKF) handles this via first-order linearisation, while the 
Unscented Kalman Filter (UKF) \cite{julier1997newextension} propagates sigma points 
through the nonlinear transformation for improved accuracy when the nonlinearity is 
significant.

In the context of multi-target tracking, the Kalman filter can not be used until 
detections have been associated with targets. To determine which detections 
correspond to which targets a common strategy involves gating target candidates 
for each detection by finding the distance between the predicted measurement for 
each target, and the actual measurement. 
After gating, and association algorithm such as the Global Nearest Neighbour 
(GNN) algorithm can create a one-to-one assignment between detections 
and targets.

\begin{figure}[htb]
    \incfig[0.5]{gnn}
    \label{fig:gnn}
    \caption[Global nearest neighbour data association.]{Global nearest neighbour 
    data association algorithm associating three targets measurement predictions 
    in red with four detections in yellow. The distance based gates are depicted 
    as ellipses around the predicted measurements. The algorithm finds a one-to-one 
    assignment between targets and detections by minimising the overall distance 
    between predicted measurements and actual measurements. In this case, the assignment
    would be T1-D1, T2-D3 and T3 with no detection.}
\end{figure}

In dense clutter, hard assignment methods can be insufficient. Probabilistic data 
association (PDA) updates each track as a weighted combination of all gated detections, 
with weights proportional to measurement likelihood. Joint PDA (JPDA) extends this to 
the multi-target case \cite{barshalom1988tracking}. Multiple hypothesis tracking (MHT) 
defers the association decision by maintaining a tree of hypotheses across multiple 
scans \cite{reid1979mht} (see Figure \ref{fig:mht}).

\begin{figure}[htb]
    \incfig[0.8]{mht}
    \label{fig:mht}
    \caption[Multiple hypothesis tracking with two targets in clutter.]{Multiple hypothesis tracking algorithm maintaining branching tracks
    of two targets (in blue and red) with their measurements indicated by blue
    dots.
    "Clutter" detections that make the path of the targets unclear are indicated
    as triangles. The blue and red shaded circles represent several potential
    posterior estimates of each of the two targets at the current timestep,
    reflecting the ambiguity in associating the measurements to the targets
    amongst the clutter. There is also an uncertain "no detection' hypothesis
    for each target indicated by the large circles. The size of the shaded circles 
    indicate 3 standard deviations of uncertainty.}
\end{figure}

Detections that cannot be associated with any existing track initialise tentative 
tracks, which are promoted to confirmed only after receiving at least M detection
associations in the last N timesteps, where M and N are parameters of the promotion 
policy, commonly referred to as M-of-N logic 
\cite{blackman1999modern}. This prevents transient clutter returns from generating 
false confirmed tracks. Conversely, tracks accumulating consecutive missed detections 
are allowed to coast for a limited number of updates before being deleted. The 
promotion and deletion thresholds balance responsiveness to new targets against 
robustness to clutter and intermittent detection loss and is generally tuned to 
fit the operating environment. 
